\section{Discussion}\label{sec:discussion}

\paragraph{Information is not decision value.}
Test (a) asks whether dispersion ranks view errors correctly; test (b) asks whether it improves the portfolio. The US pilot shows that the first does not imply the second. Black--Litterman uses $\Omega_{ii}^{-1}$ as a weight, so the magnitudes of $\Om$ matter, not only its ordering. A signal with a modest rank correlation but a hundredfold range acts like a strong signal in the optimizer. Recalibration addresses this in principle by mapping dispersion to the error scale. In practice the map is fitted on a few hundred noisy pairs, the view error is a small part of the realized error, and the recalibrated $\Om$ improved on the raw one without beating the conventions. For any proposed source of view uncertainty, LLM-based or not, we recommend reporting both tests, together with the He--Litterman benchmark and the noise-adjusted coverage of Eq.~\eqref{eq:decomp}.

\paragraph{The views are the binding constraint.}
The framework can at best allocate trust among views; it cannot create predictive content. With prompt A the views are close to a function of the last two weeks of returns, and short-term continuation was absent in 2025, in line with the reversal literature \citep{jegadeesh1990evidence,lehmann1990fads}. Under these conditions an $\Om$ that concentrates weight on a subset of views can only add variance. The constant and He--Litterman $\Om$, which weight the views more evenly, lose least, and even the shuffled $\Om$, which breaks the link between confidence and view, outperforms the correctly assigned one. Whether the dispersion mapping has value when the views do carry information is the open question. Inputs such as company news or filings could supply that information, as in \citet{lopez2023can}. They are outside the scope of this paper, which keeps the earlier version's information set fixed so that the re-examination isolates the effects of data, models and evaluation.

\paragraph{Why the earlier result did not survive.}
The earlier version used the 50 largest S\&P~500 stocks as of the end of its sample (March 2025), models including some whose developers do not document a training cutoff, and a single $\tau$ tuned on a three-month validation window and then held fixed. Each of these can produce outperformance without skill in the views: an end-of-sample universe tilts toward stocks that grew during the test period, an undocumented cutoff cannot rule out recall of realized returns, and a $\tau$ chosen on three months of data can select a view weight that worked by chance and is never revisited. In a preliminary re-run on the earlier data with a corrected pipeline, the reported gains already disappeared \todo{report the stage-0 re-run (BL-Llama Sharpe 1.23 $\rightarrow$ $-0.12$) in the appendix with its exact design}. \todo{Quantify each correction separately: fixed vs.\ point-in-time universe; in-sample vs.\ walk-forward $\tau$.}

\paragraph{Cost.}
The dispersion signal saturates at $N=5$--10 queries per stock (Table~\ref{tab:calib}). The US collection with $N=20$ took 4.5 GPU-hours on one H100 for 1{,}750 stock-dates; $N=10$ halves this. Cost is not what limits the framework.

\paragraph{Limitations.}
The test window is short by construction: models must have documented cutoffs before it and data end in December 2025, leaving 16 months per market and 26 evaluated decisions. Bootstrap intervals are correspondingly wide, and the walk-forward $\tau$ is unstable early on. Pooling four markets is our response \pending{KR, JP, DE}. This draft reports one model. A reasoning model's dispersion includes variation in its reasoning path, which may differ from that of a non-reasoning model \pending{Gemma 3 27B, Llama 3.3 70B, gpt-oss-120b}. Dispersion also depends on the sampling temperature; we use 1.0 throughout \pending{temperature 0.5 / 1.0 / 1.5 for Gemma 3 27B}. Finally, documented cutoffs bound but do not rule out exposure to later information, for example through retrieval or post-training data; we use the models offline without retrieval.
